\begin{table}[htbp]
\centering
\scriptsize
\caption{Link Between Information Loss and Disaggregation Gains}
\label{tab:cancellation_forecast_link}
\begin{tabular}{lcccccc}
\toprule
Group & $n$ dates & Pearson $r$ & HAC slope & $p$-value & Mean gain low-loss q & Mean gain high-loss q \\
\midrule
All models & 27 & 0.131 & 836,839.5 & 0.379 & -90,672.7 & 116,981.1 \\
Regime-switching & 27 & 0.273 & 547,689.8 & 0.021 & -3,778.3 & 36,646.9 \\
Other models & 27 & 0.106 & 952,499.4 & 0.486 & -125,430.4 & 149,114.8 \\
\bottomrule
\end{tabular}
\begin{tablenotes}
\footnotesize
\item \textit{Notes:} Rows use the aligned forecast panel merged with the cancellation series on common dates (2023:01--2025:03; 27 monthly observations). For each group, the dependent variable is the date-level mean squared-loss differential, $\ell_{\mathrm{agg}} - \ell_{\mathrm{disagg}}$, so positive values indicate that disaggregation lowers forecast loss. The HAC slope comes from a regression of that date-level mean loss gain on information-loss potential; low/high quartiles split dates by the 25th and 75th percentiles of information-loss potential within the same sample.
\end{tablenotes}
\end{table}
